\startenvironment fduthesis_env

% ====================================================================
%  复旦大学学位论文 ConTeXt 模板 —— 环境定义
%  支持本科 / 硕士 / 博士学位论文
%  遵循《复旦大学博士、硕士学位论文规范（2024 年 10 月修订版）》
% ====================================================================

% --------------------------------------------------------------------
%  一、论文信息变量（用户在主文件中通过 \setvariables 覆盖默认值）
% --------------------------------------------------------------------
\setvariables [thesis] [
  type=master,           % doctor / master / bachelor
  degree=academic,       % academic / professional
  secret-level=none,     % none / 秘密 / 机密 / 绝密
  secret-year={},
  school-id={10246},
  student-id={},
  title={},
  title-en={},
  author={},
  author-en={},
  department={},
  major={},
  supervisor={},
  date={},
  keywords={},
  keywords-en={},
  clc={},
  instructors={},
]

% --------------------------------------------------------------------
%  二、纸张与页面版式
%  A4 纸，上下 2.54 cm、左右 3.18 cm（Word 默认值）
% --------------------------------------------------------------------
\setuppapersize[A4][A4]

\setuplayout[
  backspace=3.18cm,
  width=middle,
  topspace=2.54cm,
  header=1.2cm,
  headerdistance=0.15cm,
  footer=1.0cm,
  footerdistance=0.4cm,
  height=middle,
  cutspace=3.18cm,
  bottomspace=2.54cm,
]

\setuppagenumbering[alternative=doublesided,location=]

% --------------------------------------------------------------------
%  三、字体配置
%  西文 Times New Roman（TeX Gyre Termes 替代）
%  中文：思源宋体 CN（Source Han Serif CN），思源黑体 CN（Source Han Sans CN）
%  正文 12pt（小四号）
%
%  字体配置说明：
%    默认使用 Source Han Serif CN / Source Han Sans CN（Adobe 思源宋体/黑体）
%    也可替换为 Noto Serif SC, SimSun 等，更换后运行：
%    mtxrun --script fonts --reload
% --------------------------------------------------------------------

\definefontfamily [mainface] [rm] [Source Han Serif CN]
\definefontfamily [mainface] [ss] [Source Han Sans CN]
\definefontfamily [mainface] [tt] [Source Han Sans CN]
\definefontfamily [mainface] [mm] [Latin Modern Math]

\setupbodyfont[mainface, 12pt]
\setscript[hanzi]

% 中文字号命令
\define\chuhao    {\switchtobodyfont[42pt]}
\define\xiaochuhao{\switchtobodyfont[36pt]}
\define\yihao     {\switchtobodyfont[26pt]}
\define\xiaoyihao {\switchtobodyfont[24pt]}
\define\erhao     {\switchtobodyfont[22pt]}
\define\xiaoerhao {\switchtobodyfont[18pt]}
\define\sanhao    {\switchtobodyfont[16pt]}
\define\xiaosanhao{\switchtobodyfont[15pt]}
\define\sihao     {\switchtobodyfont[14pt]}
\define\xiaosihao {\switchtobodyfont[12pt]}
\define\wuhao     {\switchtobodyfont[10.5pt]}
\define\xiaowuhao {\switchtobodyfont[9pt]}
\define\liuhao    {\switchtobodyfont[7.5pt]}

% --------------------------------------------------------------------
%  四、行距与段落
%  规范：行距 20 磅，首行缩进 2 字符
% --------------------------------------------------------------------
\setupinterlinespace[line=20pt]
\setupindenting[first, always, 2em]
\setupalign[hanging]
\setuptolerance[verytolerant, stretch]

% --------------------------------------------------------------------
%  五、中文本地化标签
% --------------------------------------------------------------------
\setuplabeltext  [en] [chapter={第\;,\;章}]
\setuplabeltext  [en] [figure={图\;}]
\setuplabeltext  [en] [table={表\;}]
\setupheadtext   [en] [content={目\quad 录}]
\setupheadtext   [en] [pubs={参考文献}]

% --------------------------------------------------------------------
%  六、章节标题样式
%  章标题：三号（16pt）黑体加粗，居中
%  节标题：四号（14pt）黑体加粗
%  小节标题：小四号（12pt）黑体加粗
% --------------------------------------------------------------------
\setupheads[indentnext=yes]

\setuphead [title] [
  style=\bf,
  align=center,
  before={\blank[40pt]\switchtobodyfont[16pt]},
  after={\blank[20pt]\switchtobodyfont[12pt]},
  header=empty,
  incrementnumber=list,
]

\setuphead [chapter] [
  style=\bf,
  align=center,
  before={\page\blank[40pt]\switchtobodyfont[16pt]},
  after={\blank[20pt]\switchtobodyfont[12pt]},
  header=empty,
  number=yes,
]

\setuphead [section] [
  style=\bf,
  before={\blank[20pt]\switchtobodyfont[14pt]},
  after={\blank[10pt]\switchtobodyfont[12pt]},
]

\setuphead [subsection] [
  style=\bf,
  before={\blank[10pt]},
  after={\blank[10pt]},
]

\setuphead [subsubsection] [
  style={\bf},
  before={\blank[10pt]},
  after={\blank[10pt]},
]

\setuphead[section]     [sectionsegments=chapter:section]
\setuphead[subsection]  [sectionsegments=chapter:subsection]

% --------------------------------------------------------------------
%  七、页眉页脚
%  参考 fduthesis (stone-zeng) 模板：
%    正文：偶数页左页眉=章标题，奇数页右页眉=节标题
%    前置/后置部分：页眉显示对应标题
%    页脚居中显示页码，页眉下方有横线
% --------------------------------------------------------------------

\define\FDUHeaderRule{%
  \nointerlineskip
  \vskip3pt
  \blackrule[width=\makeupwidth, height=0.4pt]%
}

\startsetups FrontHeaderFooter
  \setupheadertexts[text]
    [{\tfx\getmarking[title]}]
    [{\tfx\getmarking[title]}]
    []
    []
  \setupheader[text][after=\FDUHeaderRule]
  \setupfootertexts[text]
    [\midaligned{\tfx\userpagenumber}]
    []
    []
    [\midaligned{\tfx\userpagenumber}]
\stopsetups

\startsetups BodyHeaderFooter
  \setupheadertexts[text]
    [{\tfx 第\headnumber[chapter]章\kern1em\getmarking[chapter]}]
    [{\tfx\headnumber[section]\kern1em\getmarking[section]}]
    []
    []
  \setupheader[text][after=\FDUHeaderRule]
  \setupfootertexts[text]
    [\midaligned{\tfx\userpagenumber}]
    []
    []
    [\midaligned{\tfx\userpagenumber}]
\stopsetups

\startsetups BackmatterHeaderFooter
  \setupheadertexts[text]
    [{\tfx\getmarking[title]}]
    [{\tfx\getmarking[title]}]
    []
    []
  \setupheader[text][after=\FDUHeaderRule]
  \setupfootertexts[text]
    [\midaligned{\tfx\userpagenumber}]
    []
    []
    [\midaligned{\tfx\userpagenumber}]
\stopsetups

\startsetups EmptyHeaderFooter
  \setupheadertexts[text][][][][]
  \setupheader[text][after=]
  \setupfootertexts[text][][][][]
\stopsetups

% --------------------------------------------------------------------
%  八、图表设置
% --------------------------------------------------------------------
\setupexternalfigures[directory={., ./figures}]

\setupcaptions[
  style=\wuhao,
  headstyle=\bf,
  align=center,
  way=bychapter,
]

\setupcaption[figure][
  prefix=yes,
  prefixsegments=chapter,
]

\setupcaption[table][
  prefix=yes,
  prefixsegments=chapter,
  location=top,
]

% --------------------------------------------------------------------
%  代码排版（防止等宽字体内容溢出页面）
% --------------------------------------------------------------------
\setuptyping[
  style=\wuhao,
  before={\blank[6pt]},
  after={\blank[6pt]},
]

% --------------------------------------------------------------------
%  九、脚注
%  规范：硕士篇末注，博士章末注；也可脚注
% --------------------------------------------------------------------
\setupnote[footnote][
  way=bypage,
  numberconversion=n,
]
\setupnotation[footnote][
  way=bypage,
]

% --------------------------------------------------------------------
%  十、书签与交叉引用
% --------------------------------------------------------------------
\setupinteraction[
  state=start,
  style=,
  color=,
  contrastcolor=,
]
\setupinteractionscreen[option=bookmark]
\placebookmarks[title, chapter, section][chapter]

% --------------------------------------------------------------------
%  十一、数学公式编号（按章编号，如 (2.1)）
% --------------------------------------------------------------------
\setupformulas[way=bychapter]
\setupnumber[formula][way=bychapter]

% --------------------------------------------------------------------
%  十二、颜色
% --------------------------------------------------------------------
\usecolors[svg]

% --------------------------------------------------------------------
%  十三、目录格式
% --------------------------------------------------------------------
\setupcombinedlist[content][
  list={title,chapter,section,subsection},
  alternative=c,
]

\setuplist[title][
  style=\ss\bf,
  before={\blank[10pt]},
  after={\blank[10pt]},
  margin=0em,
]

\setuplist[chapter][
  style=\ss\bf,
  before={\blank[10pt]},
  margin=0em,
  width=4em,
]

\setuplist[section][
  margin=2em,
  width=3em,
]

\setuplist[subsection][
  margin=4.5em,
  width=3.5em,
]

% --------------------------------------------------------------------
%  十四、封面辅助宏
% --------------------------------------------------------------------

% 论文类型中文标签
\define\FDUThesisTypeLabel{%
  \doiftext{\getvariable{thesis}{type}}{%
    \doif{\getvariable{thesis}{type}}{doctor}{博\;士\;学\;位\;论\;文}%
    \doif{\getvariable{thesis}{type}}{master}{硕\;士\;学\;位\;论\;文}%
    \doif{\getvariable{thesis}{type}}{bachelor}{本\;科\;毕\;业\;论\;文}%
  }%
}

% 学位类型标签（本科不显示）
\define\FDUDegreeLabel{%
  \doifnot{\getvariable{thesis}{type}}{bachelor}{%
    \doif{\getvariable{thesis}{degree}}{academic}{（学术学位）}%
    \doif{\getvariable{thesis}{degree}}{professional}{（专业学位）}%
  }%
}

% 封面信息行（标签 4em + 冒号 + 下划线值域 12em，总宽 17em 保证居中对齐）
\define[2]\FDUCoverInfoLine{%
  \dontleavehmode
  \hbox to 4em{#1\hfil}：%
  \rlap{\lower5pt\hbox{\blackrule[width=12em, height=0.4pt]}}%
  \hbox to 12em{\hfil#2\hfil}%
  \par
  \blank[8pt]
}

% 封面信息行（宽标签，用于"专业学位类别（领域）"，标签 9em + 冒号 + 值域 7em = 17em）
\define[2]\FDUCoverInfoLineWide{%
  \dontleavehmode
  \hbox to 9em{#1\hfil}：%
  \rlap{\lower5pt\hbox{\blackrule[width=7em, height=0.4pt]}}%
  \hbox to 7em{\hfil#2\hfil}%
  \par
  \blank[8pt]
}

% ── 封一：论文封面 ──
\define\makecoveri{%
  \setups{EmptyHeaderFooter}
  \setuppagenumber[number=0]
  \page
  \setupheader[state=empty]
  \setupfooter[state=empty]
  \startalignment[center]
    %% 密级（如有，显示在右上角）
    \doifnot{\getvariable{thesis}{secret-level}}{none}{%
      \startalignment[flushright]
        {\bf 密级：\getvariable{thesis}{secret-level}\getvariable{thesis}{secret-year}}%
      \stopalignment
      \blank[12pt]
    }
    %% 学校代码与学号（右对齐，带下划线）
    \startalignment[flushright]
      \dontleavehmode
      \hbox to 4em{学校代码\hfil}：%
      \rlap{\lower5pt\hbox{\blackrule[width=8em, height=0.4pt]}}%
      \hbox to 8em{\hfil\getvariable{thesis}{school-id}\hfil}\par
      \blank[4pt]
      \dontleavehmode
      \hbox to 4em{学\hfil 号}：%
      \rlap{\lower5pt\hbox{\blackrule[width=8em, height=0.4pt]}}%
      \hbox to 8em{\hfil\getvariable{thesis}{student-id}\hfil}\par
    \stopalignment
    \blank[60pt]
    %% 校名（使用毛笔字体图片，居中）
    \midaligned{\externalfigure[fudan-name.png][width=10cm]}
    \blank[20pt]
    %% 论文类型
    {\xiaoerhao\bf \FDUThesisTypeLabel}
    \blank[10pt]
    %% 学位类型
    {\sihao \FDUDegreeLabel}
    \blank[40pt]
    %% 中文标题
    {\xiaoerhao\bf \getvariable{thesis}{title}}
    \blank[20pt]
    %% 英文标题
    {\sihao\bf \getvariable{thesis}{title-en}}
    \blank[60pt]
    %% 信息栏
    {\sihao
      \FDUCoverInfoLine{院\hfil 系}{\getvariable{thesis}{department}}
      \doif{\getvariable{thesis}{degree}}{professional}{%
        \FDUCoverInfoLineWide{专业学位类别（领域）}{\getvariable{thesis}{major}}
      }
      \doifnot{\getvariable{thesis}{degree}}{professional}{%
        \FDUCoverInfoLine{专\hfil 业}{\getvariable{thesis}{major}}
      }
      \FDUCoverInfoLine{姓\hfil 名}{\getvariable{thesis}{author}}
      \FDUCoverInfoLine{指导教师}{\getvariable{thesis}{supervisor}}
      \FDUCoverInfoLine{完成日期}{\getvariable{thesis}{date}}
    }
  \stopalignment
  \page
}

% ── 扉页：指导小组成员 ──
\define\makecoverii{%
  \setups{EmptyHeaderFooter}
  \page
  \setupheader[state=empty]
  \setupfooter[state=empty]
  \startalignment[center]
    \blank[80pt]
    {\sanhao\bf 指导小组成员}
    \blank[60pt]
    {\xiaosihao
      \getvariable{thesis}{instructors}
    }
  \stopalignment
  \page
}

% ── 封三：独创性声明与授权声明 ──
\define\makecoveriii{%
  \setups{EmptyHeaderFooter}
  \page
  \setupheader[state=empty]
  \setupfooter[state=empty]
  \blank[40pt]
  \startalignment[center]
    {\sanhao\bf 复旦大学}
    \par
    {\sanhao\bf 学位论文独创性声明}
  \stopalignment
  \blank[40pt]
  \setupindenting[first, always, 2em]
  \indentation
  本人郑重声明：所呈交的学位论文，是本人在导师的指导下，
  独立进行研究工作所取得的成果。
  论文中除特别标注的内容外，不包含任何其他个人或机构已经发表或撰写过的研究成果。
  对本研究做出重要贡献的个人和集体，
  均已在论文中作了明确的声明并表示了谢意。
  本声明的法律结果由本人承担。
  \blank[40pt]
  \startalignment[flushright]
    作者签名：\underbar{\hskip6em}\hskip2em 日期：\underbar{\hskip6em}
  \stopalignment
  \blank[80pt]
  \startalignment[center]
    {\sanhao\bf 复旦大学}
    \par
    {\sanhao\bf 学位论文使用授权声明}
  \stopalignment
  \blank[40pt]
  \indentation
  本人完全了解复旦大学有关收藏和利用博士、硕士学位论文的规定，
  即：学校有权收藏、使用并向国家有关部门或机构送交论文的印刷本和电子版本；
  允许论文被查阅和借阅；
  学校可以公布论文的全部或部分内容，
  可以采用影印、缩印或其它复制手段保存论文。
  涉密学位论文在解密后遵守此规定。
  \blank[40pt]
  \startalignment[flushright]
    作者签名：\underbar{\hskip6em}\hskip1em 导师签名：\underbar{\hskip6em}\hskip1em 日期：\underbar{\hskip6em}
  \stopalignment
  \page
}

% --------------------------------------------------------------------
%  十五、摘要环境
% --------------------------------------------------------------------

\definestartstop[abstractcn][
  before={%
    \title{摘\quad 要}
    \setupindenting[first, always, 2em]
  },
  after={%
    \blank[20pt]
    \noindentation
    {\bf 关键词：}\getvariable{thesis}{keywords}
    \doiftext{\getvariable{thesis}{clc}}{%
      \par\noindentation
      {\bf 中图分类号：}\getvariable{thesis}{clc}
    }
    \page
  },
]

\definestartstop[abstracten][
  before={%
    \title{Abstract}
    \setupindenting[first, always, 2em]
  },
  after={%
    \blank[20pt]
    \noindentation
    {\bf Keywords: }\getvariable{thesis}{keywords-en}
    \doiftext{\getvariable{thesis}{clc}}{%
      \par\noindentation
      {\bf CLC: }\getvariable{thesis}{clc}
    }
    \page
  },
]

% --------------------------------------------------------------------
%  十六、参考文献
%  规范：GB/T 7714—2015（ConTeXt 内置最接近的为 apa 风格）
% --------------------------------------------------------------------
\setupbtx[default:list][alternative=apa]

% --------------------------------------------------------------------
%  十七、附录
% --------------------------------------------------------------------
\definehead[appchapter][chapter]
\setuphead[appchapter][
  conversion=Character,
  style={\ss\sanhao\bf},
  align=center,
  before={\page\blank[40pt]},
  after={\blank[20pt]},
]

\stopenvironment
